\documentclass{umk-matematika}

\begin{document}

\maketitlepage
    {Практическая работа №22}
    {Усечённая пирамида. Объём}
    {}

\section{Фундамент (1--4 балла)}

\textbf{Задача 1 (1 балл).} В правильной четырёхугольной усечённой пирамиде стороны оснований $6$ см и $4$ см, высота $9$ см. Найдите объём.

\textbf{Задача 2 (2 балла).} Усечённая пирамида имеет площади оснований $25$ см² и $9$ см², высота $6$ см. Найдите объём.

\textbf{Задача 3 (3 балла).} В правильной треугольной усечённой пирамиде стороны оснований $6$ см и $3$ см, высота $5$ см. Найдите объём.

\textbf{Задача 4 (4 балла).} Объём усечённой пирамиды $84$ см³, высота $7$ см, площади оснований $16$ см² и $4$ см². Проверьте формулу.

\section{Стены (5--7 баллов)}

\textbf{Задача 5 (5 баллов).} В правильной четырёхугольной усечённой пирамиде стороны оснований $8$ см и $4$ см, высота $3$ см. Найдите объём.

\textbf{Задача 6 (6 баллов).} В усечённой пирамиде площади оснований $50$ см² и $18$ см². Объём равен $196$ см³. Найдите высоту.

\textbf{Задача 7 (7 баллов).} В правильной треугольной усечённой пирамиде стороны оснований $8$ см и $2$ см, высота $6$ см. Найдите объём.

\section{Кровля (8--10 баллов)}

\textbf{Задача 8 (8 баллов).} Бетонный столб имеет форму усечённой пирамиды с квадратными основаниями. Сторона нижнего основания $0{,}8$ м, верхнего $0{,}4$ м, высота $3$ м. Найдите объём бетона для $10$ таких столбов.

\textbf{Задача 9 (9 баллов).} Насыпь для дороги имеет форму усечённой пирамиды с прямоугольными основаниями. Нижнее основание $10$ м × $10$ м, верхнее $6$ м × $6$ м, высота $3$ м. Найдите объём грунта.

\textbf{Задача 10 (10 баллов).} Водонапорная башня состоит из усечённой пирамиды (нижнее основание — квадрат $6$ м × $6$ м, верхнее — $3$ м × $3$ м, высота $10$ м) и куба $3$ м × $3$ м × $3$ м сверху. Найдите общий объём.

\newpage
\section*{Ответы}

\begin{enumerate}
  \item $228$ см³
  \item $98$ см³
  \item $\frac{105\sqrt{3}}{4}$ см³
  \item данные противоречат формуле объёма усечённой пирамиды
  \item $112$ см³
  \item $6$ см
  \item $42\sqrt{3}$ см³
  \item $11{,}2$ м³
  \item $196$ м³
  \item $237$ м³
\end{enumerate}

\end{document}